\chapter{Συμπεράσματα}\label{ch:conclusions}

\section{Σύνοψη}\label{sec:concl-summary}
Η παρούσα εργασία μελέτησε το πρόβλημα της σύστασης κανόνων IoT σε ομοσπονδιακό πλαίσιο, με αφετηρία τον διαγωνισμό της Wyze και την υλοποίηση της δεύτερης θέσης. Ολόκληρη η πειραματική διαδικασία στηρίχθηκε σε ένα ενιαίο πλαίσιο προσομοίωσης, ικανό να εκτελέσει το ίδιο ακριβώς μοντέλο κάτω από διαφορετικούς τύπους ενορχήστρωσης μεταβάλλοντας μόνο τις παραμέτρους της κατανομής των δεδομένων. Πάνω σε αυτό αναπτύχθηκαν τρεις διαδοχικές εκδόσεις μοντέλου: ένα απλό baseline με GraphSAGE encoder (v1), μια έκδοση εμπλουτισμένη με τον τύπο των ακμών (v2) και, τέλος, ένα relational μοντέλο που συνδυάζει τον CompGCN encoder, τον ComplEx decoder και τη συνάρτηση σφάλματος InfoNCE με hard negatives (v3). Κάθε έκδοση αξιολογήθηκε σε τέσσερα σενάρια ενορχήστρωσης, από την πλήρως centralized εκπαίδευση έως το πλήρως cross-device περιβάλλον, και υπό δύο πρωτόκολλα αξιολόγησης που αποτυπώνουν διαφορετικές φιλοσοφίες αυστηρότητας.

\section{Βασικά συμπεράσματα}\label{sec:concl-findings}
Από την πειραματική μελέτη (\autoref{ch:experiments}) προκύπτουν τρία κύρια συμπεράσματα.

Πρώτον, η αρχιτεκτονική του μοντέλου αποφέρει σημαντικές διαφοροποιήσεις. Στο centralized σενάριο, η έκδοση v3 βελτιώνει το baseline κατά 71\% (MRR 0,4997 έναντι 0,2916). Η απομόνωση των συνιστωσών μέσω των ablation studies έδειξε ότι το κέρδος αυτό δεν προέρχεται από κάποιο μεμονωμένο στοιχείο, αλλά από τη συνεργιστική σύνθεσή τους: ο relational encoder (CompGCN) αποδείχθηκε ο καθοριστικότερος παράγοντας, ο κατευθυντικός decoder (ComplEx) υπερίσχυσε των συμμετρικών εναλλακτικών χάρη στη δυνατότητά του να αναπαριστά ασύμμετρες σχέσεις, ενώ η επιλογή της συνάρτησης σφάλματος αποδείχθηκε ότι πρέπει να ευθυγραμμίζεται με το εκάστοτε πρωτόκολλο αξιολόγησης. Καμία από αυτές τις συνιστώσες δεν επαρκούσε από μόνη της. Μόνο ο συνδυασμός τους ξεπέρασε το baseline και στα δύο πρωτόκολλα ταυτόχρονα.

Δεύτερον, η μετάβαση από την centralized στην federated εκπαίδευση έχει ένα μετρήσιμο αλλά ελεγχόμενο κόστος. Για την έκδοση v3, η επίδοση υποβαθμίζεται σταδιακά καθώς τα δεδομένα κατανέμονται σε ολοένα και περισσότερους clients (0,4997 centralized, 0,4598 cross-silo, 0,4087 pseudo-cross-device), παραμένοντας ωστόσο σαφώς ανώτερη του baseline ακόμη και στο πιο κατανεμημένο εφικτό σενάριο. Η κατανομή των δεδομένων, δηλαδή, δεν ακυρώνει την αξία του μοντέλου, απλώς εισάγει ένα προβλέψιμο τίμημα σε αντάλλαγμα μεγαλύτερης προστασίας της ιδιωτικότητας.

Τρίτον, η βέλτιστη επιλογή μοντέλου εξαρτάται από τον τύπο της ενορχήστρωσης. Στο πιο απαιτητικό σενάριο, το true cross-device περιβάλλον με έναν χρήστη ανά client και μερική συμμετοχή ανά γύρο, η κατάταξη αντιστρέφεται: το ισχυρό relational μοντέλο v3 καταρρέει (0,1981), ενώ η απλούστερη έκδοση v2 αποδεικνύεται ανθεκτικότερη (0,3791). Η παρατήρηση αυτή υπογραμμίζει ότι η πολυπλοκότητα ενός μοντέλου δεν είναι εγγενώς επωφελής: η αξία της εξαρτάται από το διαθέσιμο τοπικό context και τη σταθερότητα των federated ενημερώσεων. Όταν κάθε client διαθέτει ελάχιστη, μη-IID πληροφορία, η πολύπλοκη relational αρχιτεκτονική στερείται τη δομή που χρειάζεται για να μάθει χρήσιμες αναπαραστάσεις και υποχωρεί έναντι μιας απλούστερης λύσης.

\section{Περιορισμοί}\label{sec:concl-limitations}
Καθώς το πρόβλημα μελετήθηκε στα πλαίσια μιας διπλωματικής εργασίας, το εύρος που δύναται να καλύψει είναι πεπερασμένο, με αποτέλεσμα να έχει κάποιους περιορισμούς. Πρώτον, όλα τα πειράματα εκτελέστηκαν σε περιβάλλον προσομοίωσης και όχι σε πραγματικές edge συσκευές και ως εκ τούτου ζητήματα όπως η ετερογένεια του hardware, οι καθυστερήσεις δικτύου και οι αποσυνδέσεις clients δεν αποτυπώνονται. Δεύτερον, παρότι η ιδιωτικότητα αποτέλεσε το βασικό κίνητρο, η εργασία περιορίστηκε στη βασική στρατηγική aggregation FedAvg και δεν υλοποίησε μηχανισμούς προστασίας, όπως differential privacy ή secure aggregation. Τρίτον, η αξιολόγηση στηρίχθηκε σε ένα μόνο σύνολο δεδομένων, αυτό που παρείχε η ίδια η εταιρία που διοργάνωσε τον διαγωνισμό. Η γενίκευση των ευρημάτων σε άλλα οικοσυστήματα IoT παραμένει ανοικτό ζήτημα.

\section{Μελλοντικές επεκτάσεις}\label{sec:concl-future}
Οι παραπάνω περιορισμοί υποδεικνύουν άμεσες κατευθύνσεις για μελλοντική έρευνα. Η ενσωμάτωση μηχανισμών differential privacy και secure aggregation θα έκλεινε τον κύκλο του αρχικού κινήτρου, επιτρέποντας την ποσοτική μέτρηση της αντιστάθμισης μεταξύ ιδιωτικότητας και επιδόσεων. Η κατάρρευση της έκδοσης v3 στο true cross-device σενάριο αποτελεί κίνητρο για την διερεύνηση τεχνικών εξατομίκευσης και ανθεκτικών στρατηγικών συνάθροισης (π.χ. FedProx, FedAdam, SCAFFOLD) που αντιμετωπίζουν το client drift όταν το τοπικό context είναι ελάχιστο. Τέλος, ενδιαφέρουσες αρχιτεκτονικές επεκτάσεις περιλαμβάνουν τον διαχωρισμό των αναπαραστάσεων ανά τύπο συσκευής και κανόνα, την αξιοποίηση πλουσιότερων χαρακτηριστικών κόμβων, καθώς και τη μοντελοποίηση της χρονικής εξέλιξης των αυτοματισμών ενός χρήστη.